\section{Response-Aware and Identifiability-Adaptive Selection}
\label{sec:adaptive}

\subsection{The adaptive rule}

The third measurement of \cref{sec:diagnostics} converts directly into a policy.
Let $\bar\nu$ be the mean novelty of the candidate pool, computed on the
training split. The rule is
\begin{equation}
 \pi(A)=
 \begin{cases}
  \text{retrieve the top-}B\text{ candidates by relevance}, & \bar\nu \ge \tau,\\[2pt]
  \text{include the whole pool}, & \bar\nu < \tau,
 \end{cases}
 \label{eq:adaptive-rule}
\end{equation}
with the threshold $\tau$ chosen on held-out targets and then fixed. The rule
uses no labels at decision time, requires one covariance estimate, and replaces
a default that is harmful in the low-novelty regime. It is deliberately not a
new ranking heuristic: its purpose is to make an existing heuristic safe by
telling the practitioner when it applies. \Cref{sec:experiments} reports it on
two datasets with held-out threshold selection.

\subsection{When the decision can be learned: response-aware routing}

When the decision is estimable, the theory prescribes what to learn. The router
used throughout this paper keeps a cached view for screening and reranks the
retained candidates with a utility model trained on realised marginals. Two
details follow from \cref{sec:theory} rather than from tuning.

\paragraph{Features.} The utility model receives state features, candidate
features and \emph{response} features: the predictor's prediction before and
after adding the candidate, the mean, standard deviation, mean absolute value,
root-mean-square and maximum absolute value of the prediction change. By
\cref{prop:suff}, the full response vector is part of a sufficient
representation of the squared-loss utility conditional on the state. The
compact statistics used here are a low-dimensional approximation motivated by
that structure, not sufficient statistics in the formal sense.

\paragraph{Head.} A generic scalar head over concatenated features has to
discover the bilinear structure of \eqref{eq:prop-suff} from data. We therefore
also evaluate a \emph{bilinear response head} that emits a vector in prediction
space and scores
\begin{equation}
 \widehat\marginal_B(j\mid A)
 =
 \langle g_\theta(\phi_A(x), c_j),\, \dAj(x)\rangle
 -\|\dAj(x)\|^2-\lambda c_j ,
 \label{eq:bilinear-head}
\end{equation}
where $\phi_A(x)$ collects the state and candidate features and $g_\theta$ is a
small network. This is the minimal-capacity model consistent with the identity:
it cannot represent a score that ignores the response, it never has to learn
the curvature, and the only statistical task left is estimating the
state-dependent coefficient. For cross-entropy the exact coefficient is
replaced by the bound of \cref{prop:ce}, and we evaluate both the fixed bound
and a learned, signed coefficient.

\paragraph{Stopping.} Stopping is a separate decision from ranking. A router
that stops whenever the predicted marginal is non-positive assumes its scores
are calibrated; when the head is misspecified the rule truncates selection and
destroys value, which \cref{sec:experiments} observes on a second task family.
We therefore report both the thresholded rule and a forced-budget variant that
always spends the budget, and we treat the difference between them as a
calibration diagnostic rather than a hyperparameter.

\subsection{Operating characteristics of the rule}

A rule that must be run before a deployment is a decision procedure, so we
measured how much of the deployment has to be labelled before its verdict can be
trusted. On the six ridge-regime deployments, pilots of $n$ targets were drawn
without replacement from the 32 available (400 pilots per dataset and size) and
scored by the same rule applied to the full sample. A pilot of $24$ targets
reproduces the full-sample verdict $96\%$ of the time, with $98\%$ sensitivity
and $96\%$ specificity for the operationally critical ``route'' call; sixteen
targets give $84\%$, $86\%$ and $93\%$, and four targets already get the
\emph{direction} right almost always while agreeing on significance only
$.64$ of the time. Half the deployments are genuinely positive, so a
constant-``route'' rule would score specificity zero: the numbers are not a
base-rate artefact. In the ridge regime the pilot costs about two minutes; with
the nonlinear backbone, about half an hour.

\paragraph{Does the pilot survive a time shift?} A pilot is necessarily labelled
before the period the router will serve, so the procedure was re-run with an
added evaluation split: the router and the pool-everything baseline are scored
on the calibration period and on the test period for the same 32 targets of each
dataset, 192 targets in all. Direction transfers well --- the within-dataset
correlation between the pilot and deployment differences is $+.596$
($p<10^{-19}$), every dataset is individually significant ($+.520$ to $+.848$),
and the per-target signs agree $81\%$ of the time. The significance call
transfers less well: the dataset-level verdict agrees on four of six, against
the $96\%$ of the same-period pilot. Both disagreements keep the correct sign,
and, importantly, none of the four pilots that said ``route'' was contradicted
by a deployment on which routing hurt. The honest reading is that the procedure
is directionally robust to the time shift and its confidence is not, so the
pilot should be enlarged rather than abandoned.

\paragraph{Recognising a pilot that should not be trusted.} The pilot can be
audited against itself. Splitting it into two halves and computing the paired
interval on each gives a check whose operating characteristics are measured on
the same six deployments: at sixteen targets, pilots whose halves both resolve in
the same direction agree with the full-sample verdict $99.0\%$ of the time
against $78.1\%$ for the rest, and at twenty-four targets $99.9\%$ against
$92.3\%$. The check is high-precision and low-coverage --- strict consistency
holds for only $35\%$ of sixteen-target pilots and $39\%$ of twenty-four-target
ones --- so it certifies a verdict rather than producing one. Requiring merely
that the halves share a sign is worth almost nothing ($85.4\%$ to $84.8\%$ at
sixteen targets). The refined procedure is therefore: split the pilot, act when
both halves resolve in the same direction, and enlarge the pilot otherwise,
rather than falling back on a cheaper proxy. One caveat is stated rather than
hidden: requiring both halves to resolve selects for large effects, which are
also the ones that resolve on the full sample, so the check is partly a
signal-strength filter. A genuinely marginal deployment yields an inconclusive
check, and the correct response there is more labels.

The procedure that follows from this paper is then: label a pilot of roughly two
dozen targets, run the protocol unchanged on it --- screening, budget, objective,
and both the router and the pool-everything baseline --- compute the paired
target-level interval for router minus pooling, and route only if it excludes
zero above. The cheaper substitutes do not work, and we report each as a
negative: a fitted predictability probe is indistinguishable from chance
(\cref{sec:experiments}), structural headroom does not order the benefit in
either regime, and the cached screen is uncorrelated with the true marginal
precisely when the decision is close.

\subsection{What the theory does not claim}

The bounds of \cref{sec:theory} are conditional on the state and require the
loss gradient at the current prediction. At decision time the outcome is
unobserved, so the router estimates the expectation of the first-order term
from training episodes. The theory therefore explains why response features are
the right observation and how large the second-order error of a response-based
surrogate is; it is not a label-free guarantee, and nothing above asserts that a
learned router attains the bound.
